% =============================================================================
%  扉页、摘要、译者说明
% =============================================================================
\begin{titlepage}
\centering
\vspace*{2.6cm}
{\Huge\bfseries 一般线性群的\\[0.25em] 局部 Langlands 对应\par}
\vspace{1.0em}
{\large\itshape Notes on Local Langlands for $\GL_n$\par}
\vspace{2.2em}
{\large Zijian Yao\par}
\vspace{0.4em}
{\normalsize 哈佛大学学习讨论班讲义，2017 年秋\par}
\vspace{2.6em}
{\large 中译注本\par}
\vspace{0.6em}
{\normalsize 数译 · Langlands 导读系列 · SY02-43\par}
\vfill
{\small
\begin{tabular}{@{}rl@{}}
底本 & Zijian Yao, \emph{Notes on Local Langlands for $\GL_n$},\\
     & 哈佛学习讨论班讲义（2017 年秋），36 页（pdfTeX 排版）\\
译者 & 数译（DSH agent：lead）\\
版本 & v1.0，2026 年 11 月 27 日\\
体例 & 分节编号 1.1.1、2.3.1 等照底本；\\
     & 定理、引理、推论共用一套按节重排的编号；\\
     & 底本未写完的小节照实标注、不代拟；\\
     & 译注（橙色框）与「本章要点」（蓝色框）为译者所加\\
\end{tabular}}
\end{titlepage}

\section*{摘要}
\addcontentsline{toc}{section}{摘要}

这是 2017 年秋季哈佛一次学习讨论班的极粗糙笔记，由作者最初几次报告扩充而来。
由于报告人与听众中都没有专家，笔记里多半有我们尚未察觉的错误，阅读时请谨慎；
若发现关键性的误解，作者希望读者以电子邮件告知。

全书分六节。第 1 节整理 Weil--Deligne 表示与 $l$-进表示，以及两者之间的
$\mathrm{WD}$ 函子。第 2 节从光滑表示、容许表示、Hecke 模讲到超尖点表示与
Jacquet 函子，最后陈述 $\GL_n$ 的局部 Langlands 对应。第 3 节照 Bernstein--Zelevinsky
的分类，把 $\GL_n$ 的容许表示组织成 Zelevinsky 段的乘积，并讨论平方可积、
缓增、一般（generic）表示与 $K$-球面表示。第 4 节讨论惯性 Langlands 对应与
Bernstein 分解、类型论。第 5、6 节在底本中只有标题，正文告缺。

\begin{yizhu}[关于「粗糙笔记」的说明]
底本扉页自述 these notes are extremely rough，且多处明确写着
「to be completed in second draft」。译文对这类空缺一律照实保留（标注
\emph{底本此处未写}），不替作者补写内容；凡译者补出的背景，一律放进橙色译注框。
\end{yizhu}

\section*{译者说明}
\addcontentsline{toc}{section}{译者说明}

\paragraph{底本与体例}
底本为 Zijian Yao 2017 年秋在哈佛学习讨论班的讲义
\emph{Notes on Local Langlands for $\GL_n$}，36 页，无正式出版信息，
作者自印 pdfTeX 稿（文件内嵌日期 2019-10-11）。底本的特征有三点需要在译文里交代：

\begin{enumerate}[leftmargin=2em]
\item \textbf{分层写作}。作者按「第一稿写骨架、第二稿补内容」推进，因此第 1、2 节基本完整，
      第 3 节已开始出现空小节（3.3.2、3.4、3.6.2），第 5、6 节只剩标题。
      译文保留全部小节编号与标题，空缺处标注\emph{底本此处未写}。
\item \textbf{大量「注」「例」不编号}。底本的定义、注、例一律不编号，
      编号只给定理、引理、推论、命题，且四者共用一个按节重排的计数器
      （例如 1.3 是定理，1.4 是推论）。译文照此办理。
\item \textbf{记号不统一处照录}。底本在 1.1.4 用 $V(-1)$ 表示形式符号意义下的扭转，
      又用 $\epsilon: G_K \to L^\times$ 表示对应的非分歧特征；正文两处都保留，
      但在译注里点明这是「形式符号 + 具体实现」两层意思。
\end{enumerate}

\paragraph{术语}
译文术语以本库已冻结的译名为准，并与香蕉空间（bananaspace.org）现有词条交叉核对：
「局部 Langlands 对应」「惯性群」「尖点表示」「导子」「非分歧」「分歧」「驯分歧」
在该站有页面或导航框用例，直接沿用；「超尖点表示」「容许表示」「Weil--Deligne 表示」
在该站无专条，沿本库既有译名（超尖点、容许的、Weil--Deligne 表示）。
逐条对照见 \texttt{资料卡/04-术语文献勘误.md}。

\paragraph{译注的分工}
橙色框只做三件事：补底本省略的推理、点明记号冲突、交代该概念在中文文献里的别名。
凡涉及数学内容的判断，都以「底本如此」与「译者按」两类措辞区分，
不与原文陈述混同。

\paragraph{致谢与免责}
底本作者明言笔记可能有错。译文对底本中明显的笔误（如把 $\GL_n$ 写成 $\GL_2$、
把 $\sigma$ 写成 $\psi$）在译注中列出，但不改动正文公式；读者若发现译文自身的问题，
请以底本为准。
